\documentclass[12pt]{article}
\usepackage[margin=1in]{geometry}
\usepackage{hyperref}
\usepackage{graphicx}
\usepackage{enumitem}

\title{SciCalc Next CI/CD Report}
\author{Pradyun Devarakonda}
\date{\today}

\begin{document}
\maketitle

\section{Introduction}
This report documents the implementation of the SciCalc Next scientific calculator application and its accompanying DevOps pipeline. It summarises the project scope, tooling, and the automation built to support repeatable delivery.

\section{Setup and Architecture}
\begin{itemize}[leftmargin=*]
  \item Frontend built with Next.js 15 (App Router) and TypeScript.
  \item Domain logic implemented in a dedicated \texttt{src/lib} module for reuse across API routes and unit tests.
  \item Client UI delivered through a revitalised calculator workspace with expression parsing, quick-operation toggles, and clipboard support.
  \item Visual styling leverages a glassmorphic theme, responsive panel/sidebar layout, and semantic labelling for improved accessibility.
\end{itemize}

\section{CI/CD Implementation}
The Jenkins pipeline installs dependencies, runs quality gates (linting and unit tests), executes a production build, and publishes a Docker image to Docker Hub. Credentials are injected securely via Jenkins.

\section{Docker and Ansible Deployment}
A multi-stage Dockerfile produces a lean runtime image based on \texttt{node:20-alpine}. Ansible playbooks pull the published image and ensure the container remains running with an exposed port of 3000.

\section{Results and Screenshots}
The redesigned interface surfaces contextual guidance, responsive layouts, and clipboard-enabled results. Captured UI previews are tracked under \texttt{../screenshots/} for auditability:
\begin{itemize}[leftmargin=*]
  \item Screenshot 2025-09-30 at 21.03.48 --- dark theme overview of the calculator shell.
  \item Screenshot 2025-09-30 at 21.10.27 --- expression parsing and quick action controls in use.
  \item Screenshot 2025-09-30 at 21.11.19 --- factorial input validation state.
  \item Screenshot 2025-09-30 at 21.50.52 --- result panel with copy feedback messaging.
\end{itemize}

\section{Learnings and Improvements}
Expression parsing introduced edge cases around negative factorials and whitespace that are now guarded inside the UI before API submission. Future iterations should add automated UI tests to cover copy-to-clipboard fallbacks and broaden accessibility audits.

\end{document}
